\section{Experimental Setup}\label{sec:exp}

\subsection{One run of the study}

We train a Transformer on the five years 2016--2020 and deploy it on
2022. Now the license on the 2020 data expires. We retrain the same
architecture, with the same hyperparameters and the same seed, on the
remaining four years. The original model is the reference of Section~\ref{sec:method},
and the retrained one is the oracle. This one pair answers every
question in the paper. Its fit gap on 2020 is the memorization gap
$G$, its forecast difference on 2022 is the deletion effect
$\Delta$, and whether an auditor can tell the pair apart is the
audit question.

The full study repeats this comparison 3{,}200 times. One
combination of architecture, training window, deleted year, and
deletion level is a \emph{cell}, with ten seeds per cell. Five
architectures, four windows, five deletable years, three levels,
and ten seeds give 3{,}000 oracle retrains, and a placebo control
adds 200.

\subsection{Forecasting task and data}

The task is 5-day-ahead realized volatility for S\&P~500 stocks, one
cross-section per trading day: each day the model ranks stocks by how
volatile the next five days will be. This is a risk task rather than
an alpha contest; the ranking feeds volatility targeting and position
sizing. For each asset $i$ and day $t$ the model reads a 64-day
history of eleven return-derived features and predicts
$y^{\mathrm{raw}}_{i,t} = \log
\operatorname{sd}\{r_{i,t+1},\ldots,r_{i,t+5}\}$, the log
volatility of its next five daily returns, $z$-scored within each
day's cross-section and clipped to $[-3,3]$. Training minimizes
mean squared error against this target.

Daily returns are dividend- and split-adjusted simple returns for a
fixed list of 501 current S\&P~500 constituents, so delisted firms
are excluded by construction. The panel covers 2015--2025, about
480 stocks per day, and hyperparameter tuning uses an earlier
extract reaching back to 2008. The panel and its construction
scripts are released.

\subsection{Rolling windows and deletions}

There is one training design, not four: a five-year window slid
forward one year at a time. Each position of the window validates
on the following year and tests on the year after that. We call
each window position a \emph{fold}, named by its test year: the
2022 fold trains on 2016--2020, validates on 2021, and tests on
2022, and the folds run through 2025. Rolling fixed-width windows
are standard deployment practice and give every era a scheduled
exit, which makes deleting a year still inside the active window
the live compliance case.

Sliding the window is also what separates the paper's two questions.
Within one window position, the deleted year varies while the test
year is held fixed. Across window positions, the same deleted year
(2020, the crisis) is removed under four different test years. The
first isolates what was deleted, and the second isolates the market
it is judged in.

Two controls complete the design. A \emph{random-deletion placebo}
removes the same number of windows as the 2020 deletion, drawn
uniformly, and separates the effect of losing data volume from the
effect of losing that particular year. We deliberately do not top up
the training set to equalize volume, because the compliance
counterfactual of an expiry is a smaller training set.

\subsection{Models and benchmark}

We tune seven forecasting backbones and report the five that span the
fit range: the vanilla Transformer~\cite{vaswani2017attention},
TSMixer~\cite{chen2023tsmixer}, SegRNN~\cite{lin2023segrnn},
iTransformer~\cite{liu2024itransformer}, and
PatchTST~\cite{nie2023patchtst}. Two further backbones fit worse
than the weakest of these and add no further contrast, so they
appear only in the released sweep.

As an econometric benchmark we refit the heterogeneous autoregressive
(HAR) model of realized volatility~\cite{corsi2009har} under the
identical protocol: log 5-day-forward realized volatility regressed
on its trailing 1-, 5-, and 22-day counterparts by ordinary least
squares. It has four parameters, is subject to the same deletions,
and sets the line the results are organized around: a neural
forecaster that cannot beat it has no business in production. We call
the three architectures that beat it in every fold the
\emph{deployable} models. The two that never do underfit this task,
and we call them the \emph{non-deployable} models.

\subsection{Metrics, tuning, and cost}

Forecast skill is the daily cross-sectional information coefficient
(IC): each day, the Spearman rank correlation between predictions and
targets across stocks, averaged over the test year's days,
written $s_{\text{test}}$. The same quantity computed on the deleted
year's windows is written \forgetic{}. The forget-set loss
$\mathcal{L}_F$ is the mean squared error on those windows. The
governance metrics $G$, $\Delta$, and $\gapclosed$ of
Eqs.~\ref{eq:gap}--\ref{eq:kappa} are built from these. Every
comparison is paired by trading day with Newey--West standard
errors~\cite{newey1987hac}, forget-set metrics are means over ten
seeds, and the audit separation $d'$ compares the two ten-seed
populations.

Hyperparameters were tuned once, on a split that predates every
deletable year (2010--2014, validated on 2015, disjoint seeds),
selecting on validation loss over a grid of width and dropout. A second search over learning rates and structural
settings changed nothing. iTransformer and
PatchTST remain below the HAR benchmark after that dedicated search,
so their weak fit is a property of the architecture on this task, not
of the tuning budget. A single oracle retrain costs 3--18 minutes on
one NVIDIA A100, and the full campaign took about 600 A100-hours.
Full configurations are in the released protocol.
